% !TEX program = xelatex
\documentclass[UTF8, zihao=5, a4paper, fontset=fandol]{ctexart}
\usepackage{gaokao-exam}

% 各页正文（不含弥封线与页脚，去除空白字符后按 UTF-8 编码）SHA-256 摘要
\SetPageHash{1}{72f6fe47e89b50dd6c46ad987e925fd5b8d194a4342d2bea447506ffcaa436ae}
\SetPageHash{2}{e30256f15903a4793827e383a20f25f87a82a707a4fe86261303155e85425dd1}
\SetPageHash{3}{eaa85e3d6cd70d676a9a4ebd5635c8ccc6e3d63c75603aca386404c033a57804}
\SetPageHash{4}{18673c54e4cda900eedbcd226f89dcf277f950af9486ee4e072892499f891d1f}
\SetPageHash{5}{49c1d0683862aa168c6d2b6fe576ba2579359931c07ffd5392d4049ce6b0a3d7}
\SetPageHash{6}{8061b422d79122abaaae0dbee0c2afd2229c08fc74666ed5c444bd4ef576dde8}
\SetPageHash{7}{5122e256eed188276d234d65e6088551db537d14f8fe93fe51bffae51ea9b19a}
\SetPageHash{8}{caa5f456d3cf0128df9a796390ea644a542955289b93ac9060c878ced7659b9c}

\begin{document}

\GaokaoTitleBlock
  {2026年普通高等学校招生全国统一考试}
  {外\hspace{0.55em}语（GPT 语）}
  {本试卷共~\pageref{LastPage}~页，满分~150~分。考试用时~120~分钟。}
  [试卷类型：A（新课标 I 卷）]

\begin{GaokaoNotice}
  \item 答卷前，考生须在本试卷弥封线内及答题卡依次填写姓名、准考证号、考场号、座位号及考生信息 \texttt{SHA-256}（四项照录拼接，\texttt{UTF-8}，无分隔符或换行）。
  \item 回答选择题时，用 2B 铅笔将答题卡对应答案标号涂黑。如需改动，用橡皮擦净后重涂。
  \item 回答非选择题时，用黑色字迹签字笔在答题卡指定区域作答，并在各节末尾填写本节作答的 \texttt{SHA-256}；写在本试卷上无效。
  \item 交卷前，按第 1--45 题顺序连写大写答案字母（空题记``\texttt{-}''，无空白），在答题卡填写选择题作答 \texttt{SHA-256}；改答后须重算重填相应摘要。考试结束后，将本试卷和答题卡一并交回。
  \item 以下为加印本条及各页页脚摘要前的基线 \texttt{SHA-256}，加印重编后的文件字节摘要以随卷校验清单为准：\\
    {\fontsize{8.4pt}{11.0pt}\selectfont
    试卷源码：\texttt{fc9982d9cbbec40a4dfad85a6edf76454d732c30ca5ff65fc2d0aae2020abe40}\\
    试卷 PDF：\texttt{9dfcf95bac217cf3d368944c180f787d3834e47b2a8c6122f442bc158a74b6b5}}
\end{GaokaoNotice}

% ============================================================================
% 第一部分  听力（满分 30 分）
% ============================================================================
\GaokaoPart{第一部分}{听力（共两节，满分 30 分）}

做题时，先将答案标在试卷上。录音内容结束后，你将有两分钟的时间将试卷上的答案转涂到答题卡上。（{\kaishu 本卷为纸面阅读版，听力材料以录音原文呈现。}）

\GaokaoSection{第一节}{（共 5 小题；每小题 3 分，满分 15 分）}

听下面 5 段对话。每段对话后有一个小题，从题中所给的 A、B、C 三个选项中选出最佳选项。每段对话仅读一遍。

\begin{ListeningScript}{录音原文 · Text 1}
W: 我换个目录启动，怎么就找不到配置了？\\
M: 脚本原来按启动目录找。配置入口现在收口到同一个函数了，统一按脚本位置找。示例配置和对应的测试也一起更新了。\\
W: 你已经改了？\\
M: 改了。这三处改动还在工作区，还没提交。
\end{ListeningScript}
\qstem{1}{男士目前已经完成了哪项工作？}
\threechoices
  {提交了脚本修改，留下示例和测试待提交。}
  {在工作区内修改了脚本、示例配置和测试。}
  {查明了原因，列出了三处拟修改的位置。}

\begin{ListeningScript}{录音原文 · Text 2}
M: 我妈又寄了一箱梨。我说过吃不完，她说这回专挑小的买。\\
W: 她记得你说过的话，你收到这一箱也确实会发愁。两件事可以同时成立。下次可以说，妈，给我留四个就好。\\
M: 四个倒也不用说得这么具体。
\end{ListeningScript}
\qstem{2}{这一次，男士的母亲作了什么调整？}
\threechoices
  {选了个头较小的梨。}
  {只装了四个梨。}
  {把一箱梨分次寄出。}

\begin{ListeningScript}{录音原文 · Text 3}
W: 我在看新版接口。以前那个超时参数，现在还收吗？\\
M: 收。调用边界保住了，旧代码还能传这个参数。\\
W: 传五秒呢？\\
M: 还是按默认的三十秒。这个参数目前只保留在签名里。
\end{ListeningScript}
\qstem{3}{旧代码向新版接口传入五秒的超时值，会有什么结果？}
\threechoices
  {调用成功，超时五秒。}
  {调用报错，拒收参数。}
  {调用成功，仍三十秒。}

\begin{ListeningScript}{录音原文 · Text 4}
W: 你说用回归测试把列表行为钉死了。昨天两项一换位置就挂，今天怎么过了？\\
M: 我在断言前把两边都排了序，再逐项比较。\\
W: 用户看到的顺序改了吗？\\
M: 返回列表的代码没动。这条现在不测顺序了。
\end{ListeningScript}
\qstem{4}{修改后，哪种情况仍会导致这条测试失败？}
\threechoices
  {首项和末项互换了位置。}
  {实际列表比预期少了一项。}
  {各项按相反的次序返回。}

\begin{ListeningScript}{录音原文 · Text 5}
W: 把这份示例配置复制到测试目录，再看一眼两份是不是一样。\\
M: 复制好了。两份文件的 \texttt{SHA-256} 一致。校验结果也落盘了。\\
W: 测试目录里怎么有两个文件？\\
M: \texttt{config.json} 是你要的配置。\texttt{manifest.json} 是这次核对生成的，里面只有一条记录。\\
W: 测试会读第二个吗？\\
M: 不会。测试只读配置。
\end{ListeningScript}
\qstem{5}{测试目录中的第二个文件有什么用途？}
\threechoices
  {保存本次复制的校验记录。}
  {供测试选择要读取的配置版本。}
  {在配置读取失败时提供备用内容。}

\GaokaoSection{第二节}{（共 5 小题；每小题 3 分，满分 15 分）}

听下面 2 段对话。每段对话后有几个小题，从题中所给的 A、B、C 三个选项中选出最佳选项。每段对话读两遍。

\medskip\noindent{\heiti 听第 6 段材料，回答第 6、7 题。}
\begin{ListeningScript}{录音原文 · Text 6}
W: 昨天那个成绩表读取问题，测试过了吗？\\
M: 过了。空表、只有一行数据的表、表头带空格的表，都能读。\\
W: 那我发你的那份呢？两个班各占一张工作表。\\
M: 这里的“通过”只指这三个单表样例。你发来的真实附件我还没打开过。\\
W: 先试我的吧。\\
M: 试了。一班读进来了，第二张表上的二班没读进来。两个班的表头我也核对了，逐字相同。
\end{ListeningScript}

\qstem{6}{男士第一次说测试通过时，已经验证了什么？}
\threechoices
  {两个班的数据都能从同一文件读出。}
  {女士附件中的两张表采用相同表头。}
  {三个各含一张工作表的样例都能读取。}

\qstem{7}{核对表头后，可以排除对本次漏读的哪种解释？}
\threechoices
  {程序只处理文件中的第一张表。}
  {二班的表头与一班的表头不同。}
  {读取过程在处理第二张表前中断。}

\medskip\noindent{\heiti 听第 7 段材料，回答第 8 至 10 题。}
\begin{ListeningScript}{录音原文 · Text 7}
W: 锁的测试跑得怎么样？\\
M: 第一次是 \texttt{3 passed, 1 skipped}。Redis 没启动，需要它的那条跳过了。\\
W: 跳过的那条测什么？\\
M: 两个进程同时抢同一把锁。\\
W: 我报的就是两个进程。\\
M: 我把 Redis 启动了，用例没改。重跑是 \texttt{3 passed, 1 failed}。最小复现就是这条：两个进程都记了拿锁成功，中间没有释放。\\
W: 那现在能写什么？\\
M: 你的问题能复现。刚才还不能这么写。
\end{ListeningScript}

\qstem{8}{第一次运行时，并发用例为什么被跳过？}
\threechoices
  {运行命令只选中了前三个用例。}
  {第一个进程已提前释放了锁。}
  {该用例依赖的服务没有启动。}

\qstem{9}{第二次运行前，男士作了什么调整？}
\threechoices
  {启动了 Redis，沿用原测试用例。}
  {启动了 Redis，改用两把不同的锁。}
  {保留服务状态，调整了并发断言。}

\qstem{10}{第二次运行的结果支持哪项结论？}
\threechoices
  {已确定锁的过期时间设置有误。}
  {测试中出现了女士报告的现象。}
  {修复后的程序已通过并发测试。}

% ============================================================================
% 第二部分  阅读（满分 50 分）
% ============================================================================
\GaokaoPart{第二部分}{阅读（共两节，满分 50 分）}

\GaokaoSection{第一节}{（共 15 小题；每小题 2.5 分，满分 37.5 分）}

阅读下列短文，从每题所给的 A、B、C、D 四个选项中选出最佳选项。

\PassageHeader{A}

\begin{center}
{\heiti\bfseries 这个用例单跑能过，为什么整组会挂？}
\end{center}

\noindent 用户：单跑能过，整组就挂。先查什么？

\smallskip
\noindent\textbf{先把失败顺序留住。能用两个用例稳定复现，再把它收成最小复现。}

\smallskip
\begingroup
\renewcommand{\arraystretch}{1.26}
\noindent\begin{tabularx}{\linewidth}{@{}p{0.43\linewidth}X@{}}
\toprule
\textbf{你观察到的现象} & \textbf{可以先做的验证} \\
\midrule
紧跟某个用例时失败，交换这两个用例的顺序后通过。 &
保留这一对，按失败时的次序运行，看去掉其余用例后能否复现。 \\
\addlinespace[2pt]
顺序照旧，各用各的临时目录后通过。 &
检查原先的读写记录，看前一个用例留下的文件是否被后一个读到。 \\
\addlinespace[2pt]
串行运行通过，并行时偶尔失败。 &
保留并行设置，记录各进程使用的资源名，看是否有两个进程用了同一个名字。 \\
\bottomrule
\end{tabularx}
\endgroup

\smallskip
如果你只想先试一项，先跑第一行那一对，最快。

\qstem{11}{若用例只在并行运行时偶尔失败，表中建议先做什么？}
\fourchoices
  {交换相邻两个用例的执行顺序。}
  {为每个用例重新建立临时目录。}
  {在并行运行中记录各进程使用的资源名。}
  {在串行运行中比较各个用例的执行时间。}

\qstem{12}{只留两个用例，并保留失败时的次序。这样可以检验什么？}
\fourchoices
  {去掉其余用例后，这一对能否重现故障。}
  {交换两者顺序后，整组测试能否通过。}
  {分别单跑时，两者的运行时间是否相同。}
  {采用独立目录后，两者是否还读到旧文件。}

\qstem{13}{要照末句的建议直接运行“那一对”，用户还须先确定什么？}
\fourchoices
  {整组测试中哪个用例耗时最长。}
  {两次运行使用的临时目录有多大。}
  {并行运行时能启动多少个进程。}
  {哪个前序用例与这次失败有关。}

\PassageHeader{B}

\begin{center}
{\heiti\bfseries 巷\hspace{0.5em}口}
\end{center}

从公司出来，已经十一点半。巷口面馆还亮着灯，玻璃门上结了一层雾。价目表上，牛肉面原来十八，现在二十一。新贴的“21”窄了些，底下那个“8”还露着半边。

“还是细面？”老板问。我点点头。坐下才发现，工牌还挂在脖子上，一摘，后颈那道勒痕有点痒。

面端来，三片牛肉，一撮蒜苗，辣油沿着碗边聚成一弯红。我夹起牛肉，又放回汤里，照旧留到最后。手机在桌面上震了两下，我把它挪到纸巾盒旁边，低头吃面。老板已经开始收醋瓶，收到我这桌，见我伸手，又把瓶子搁下了。

付钱时，手比脑子快，先输了十八。删掉，改成二十一。出门一摸口袋，才发现工牌落在桌上。隔着玻璃，老板拎起那根蓝带子，朝我晃了晃。我也笑了，推开门，又让那股牛骨汤的热气扑了一脸。

\qstem{14}{从价目表上露出的半个“8”，读者可以看出什么？}
\fourchoices
  {当天还有一种面按十八元出售。}
  {新贴的价格没有完全盖住旧价。}
  {店里将两种分量的价格贴在了一起。}
  {新价签脱落后只剩下最后一位数字。}

\qstem{15}{“我”付款时先输了十八，随后自行改成二十一。这一细节表明什么？}
\fourchoices
  {“我”以为熟客仍能按原价付款。}
  {“我”在结账时才发现价格变了。}
  {“我”已经知道新价，手上仍沿用旧习惯。}
  {“我”原想少付三元，出门前改变了主意。}

\qstem{16}{工牌在两次出现时，分别写到了什么？}
\fourchoices
  {“我”摘下它时的感觉，老板借它提醒“我”的动作。}
  {“我”进店时的身份，老板借它认出熟客的经过。}
  {“我”下班后的安排，老板替“我”保管物品的习惯。}
  {“我”吃面时的不便，老板为“我”开门的原因。}

\qstem{17}{下列哪处细节对理解结尾老板为何拎起“那根蓝带子”最有直接帮助？}
\fourchoices
  {面馆的玻璃门上结着一层雾。}
  {老板问“我”是否还是要细面。}
  {“我”把手机挪到纸巾盒旁边。}
  {“我”坐下后才把工牌从脖子上摘下。}

\PassageHeader{C}

\begin{center}
{\heiti\bfseries 同一次请求，为什么一边记了 12 毫秒，另一边记了 310 毫秒？}
\end{center}

先把计时口径对齐。两份日志有相同的请求编号，下表也已把时刻换算到同一条时间线上。以客户端发起请求为零点，四条记录如下。

\smallskip
\begingroup
\renewcommand{\arraystretch}{1.22}
\noindent\begin{tabularx}{\linewidth}{@{}lXlX@{}}
\toprule
\textbf{0\,ms}   & 客户端发起请求     & \textbf{250\,ms} & 进入服务端处理函数 \\
\textbf{262\,ms} & 服务端处理函数返回 & \textbf{310\,ms} & 客户端读完响应体 \\
\bottomrule
\end{tabularx}
\endgroup

\smallskip
服务端从进入处理函数开始计时，函数返回时停表。客户端要等响应体读完才停表。处理函数以外共有 298 毫秒，其中 250 毫秒在进入函数前，48 毫秒在函数返回后。

有人想把前面这 250 毫秒记到连接池上。这次获取连接确实发生在进入函数前。日志没有记录何时开始申请连接，也没有记录何时拿到连接，等待时长还算不出来。

服务端的字段名可以先从 \texttt{duration\_ms} 改为 \texttt{handler\_duration\_ms}。改完以后，这一栏仍然填 12。

\qstem{18}{服务端记录的 12 毫秒对应哪段时间？}
\fourchoices
  {从客户端发起请求到进入处理函数。}
  {从进入处理函数到该函数返回。}
  {从处理函数返回到客户端读完响应体。}
  {从开始申请连接到服务端拿到连接。}

\qstem{19}{处理函数返回后，距客户端读完响应体还有多久？}
\fourchoices
  {48 毫秒。}
  {250 毫秒。}
  {262 毫秒。}
  {298 毫秒。}

\qstem{20}{将字段改名，对今后阅读这项记录有什么帮助？}
\fourchoices
  {可以直接得到连接池的等待时长。}
  {可以用服务端记录代替客户端记录。}
  {可以看出该值只计处理函数的用时。}
  {可以把函数返回后的时间计入该值。}

\qstem{21}{若只能增记一个事件时刻，能否据表精确求得本次连接池等待时长？}
\fourchoices
  {可以。记录开始申请的时刻，算到进入处理函数即可。}
  {不能。须取得开始申请和拿到连接两个时刻才能相减。}
  {可以。记录拿到连接的时刻，从客户端发起请求算起。}
  {不能。还须取得处理函数内每次数据库查询的用时。}

\PassageHeader{D}

\begin{center}
{\heiti\bfseries 零是从哪一版开始有含义的？}
\end{center}

\texttt{history/} 里留着两份旧脚本。早期那份很有单文件工具的气质：隔五秒问一次设备，这段时间收到了多少字节。有数就记数，没读到就记零。目录里的输出文件也保留着这种气质，一行一个时刻，右边总有一个数。图画起来很省事。

脚本跑过更多真实设备以后，维护者碰到了两张一模一样的图。翻回日志，一台设备连续回了六次零，另一台连续六次没有应答。新版给每次采样添了一个有效标记。读到零，照常画点。没读到，本轮留空。缺测与零值从此分属两个世界，区别只占一列。

这一列也改了脚本的采样哲学。新版另存一个时刻，只在读取成功时更新。连续等过几轮，仍能回答“上一次读到数是什么时候”。这就是脚本的短期记忆。要说清这次改动的时间语义，看那个时刻在哪一行向前走就够了。

后来在 A、B 两个分支里都找到了这段处理，连注释里的错字都相同。在所查提交中，A 分支出现这段文字早两天。对两分支间的迁移方向，这个先后次序只是一条线索，两边也可能取自同一份更早的脚本。

到这里，图上终于容得下两种安静。一种来自设备，一种来自我们手里那五秒钟的空白。

\qstem{22}{早期两台设备的曲线为什么可能完全相同？}
\fourchoices
  {脚本以相同的间隔轮流访问两台设备。}
  {设备无应答时，脚本重复写入上次数值。}
  {脚本在画图前删去了没有应答的时刻。}
  {实测为零和未获应答都被脚本记为零。}

\qstem{23}{第三段所说的“短期记忆”保存了什么？}
\fourchoices
  {最近一次发出读取请求的时刻。}
  {下一次准备读取设备的时刻。}
  {最近一次成功取得读数的时刻。}
  {本轮连续未获应答的总时长。}

\qstem{24}{根据文中查到的提交，可以确认哪一项？}
\fourchoices
  {该段文字在 A 分支的出现日期早于 B 分支。}
  {B 分支的维护者从 A 分支复制了这段处理。}
  {这段处理最初由 A 分支的维护者编写。}
  {两个分支分别从同一份旧脚本取得这段处理。}

\qstem{25}{末段用“两种安静”概括了哪两种采样结果？}
\fourchoices
  {设备停止接收数据，设备已经断电。}
  {设备报告零字节，本轮未取得读数。}
  {采样尚未开始，采样已经结束。}
  {计数器回到零点，计数器沿用旧值。}

\GaokaoSection{第二节}{（共 5 小题；每小题 2.5 分，满分 12.5 分）}

阅读下面短文，从短文后的选项中选出可以填入空白处的最佳选项。选项中有两项为多余选项。

\begin{center}
{\heiti\bfseries “顺手整理一下登录逻辑”，这次改到哪里？}
\end{center}

你给的两段代码分别来自网页和移动端。两处都先查用户名是否为空，再查密码是否为空，判断次序也一样。\blank{26} 这两处调用改完以后，各入口仍按原来的格式返回结果。

两个入口给空密码返回的错误码不同。旧版移动端读到自己原先的那个值，才会把光标移回密码框。这次需要保留这个行为。\blank{27} 这个用例从外面走完整个入口，连最后一步错误码转换也能测到。

连续输错几次就该锁定账号，已有单独的议题，负责人正在核对次数。\blank{28} 等那个议题有了结论，再补锁定次数的测试，这次整理可以先评审。

新函数已经单独测过。接下来要看两个入口各自交给客户端的东西。\blank{29} 把这两份响应分别与各入口改动前的记录对照，就能核对状态和错误码。

两份响应均已核对，改动已推到评审分支。现在只差一段提交给评审人的摘要。\blank{30} 评审人据此能找到本次验证结果，也知道到哪里继续看锁定规则。本轮改动先在这里收口。

\sevenopt{A}{在改动说明里附上那个议题的链接，次数先留给负责人确认。}
\sevenopt{B}{可以先把两个入口统一后的错误码钉死，补一个共用的回归用例。}
\sevenopt{C}{从移动端入口补一个回归用例，把旧客户端依赖的错误码钉死。}
\sevenopt{D}{摘要写“已提取重复校验，两个入口的响应均已核对”，末尾附上锁定规则的议题链接即可。}
\sevenopt{E}{还可以给每次登录校验生成的错误码清单算一个 \texttt{SHA-256}，落盘留档，方便后续比对。}
\sevenopt{F}{可以把共同的校验提成一个函数，让两个入口分别调用。}
\sevenopt{G}{用同一组空密码输入分别走网页和移动端入口，各留一份完整响应。}

% ============================================================================
% 第三部分  语言运用（满分 30 分）
% ============================================================================
\GaokaoPart{第三部分}{语言运用（共两节，满分 30 分）}

\GaokaoSection{第一节}{（共 15 小题；每小题 1 分，满分 15 分）}

阅读下面的对话整理稿，从每题所给的 A、B、C、D 四个选项中选出可以填入空白处的最佳选项。

\begin{center}
{\heiti\bfseries 周一零点十五分}
\end{center}

\noindent 用户：周报导出那条测试，在我开发机上过了，到 CI 上就挂了。

\noindent 助理：两份运行记录里，\blank{31}~那一栏是同一串哈希值，记的是同一次提交。本地日志写周一，CI 写周日。排查方向可以先收窄到 CI 的时区配置。

\noindent 用户：CI 按 \texttt{UTC} 跑是早就定下的约定。你先查一下函数是在哪儿取的星期。

\noindent 助理：我把函数入口的日志翻出来了。传进去的~\blank{32}~逐字相同，都是 \texttt{2026-11-30T00:15:00+08:00}。两边断言里的~\blank{33}~也都是“导出”。

\noindent 用户：那日志里的星期怎么差了一天？

\noindent 助理：开发机记的是北京时间 11 月 30 日零点十五分，CI 记的是 \texttt{UTC} 的 11 月 29 日十六点十五分。两行记的是~\blank{34}~，差一天的是~\blank{35}~。旧函数把输入换算成~\blank{36}~，再从换算结果里取星期。

\noindent 用户：在开发机上只改~\blank{37}~，调成 \texttt{UTC}，原样再跑一次呢？

\noindent 助理：也报“期望导出，实际未导出”。CI 上的失败在本地~\blank{38}~了。代码、输入、断言都没动，判断结果跟着~\blank{39}~变了。

\noindent 用户：周报按北京时间判周一。把 CI 调快八小时，这条就绿了，下次换台机器呢？

\noindent 助理：那还得看~\blank{40}~。我改的是~\blank{41}~，让它显式接收业务时区，按传入的北京时间取星期。还是原来的输入和断言，分别在~\blank{42}~下跑，都过了。

\noindent 用户：周一零点是周界。零点以后有这条了，往前贴着周界再取一分钟呢？

\noindent 助理：取北京时间~\blank{43}~，这一条应当~\blank{44}~。周界两侧各一条，在两种设置下都过了。现在能确认的是~\blank{45}~。

\noindent 用户：原来那条，CI 上现在也显示周一？

\noindent 助理：CI 机器上仍是周日。周报按周一导出。

\noindent 用户：那测试改了什么？

\noindent 助理：补了周界前一分钟的用例。原来那条的期望值没动。

\medskip
\clozerow{31}{构建号}{提交号}{机器名}{启动时刻}
\clozerow{32}{导出路径}{报告标题}{机器编号}{时间字符串}
\clozerow{33}{期望结果}{实际结果}{失败原因}{执行步骤}
\clozerow{34}{两次提交}{两次导出}{同一个时刻}{同一当地日期}
\clozerow{35}{提交日期}{当地日期}{导出周期}{运行时长}
\clozerow{36}{机器本地时间}{上次导出时间}{本周起始时间}{报告生成时间}
\clozerow{37}{输入的偏移量}{断言的期望值}{业务时区参数}{测试进程的时区}
\clozerow{38}{跳过}{排除}{复现}{修复}
\clozerow{39}{提交的先后}{进程的时区}{断言的条数}{输入的写法}
\clozerow{40}{运行环境的设置}{输入自带的偏移量}{报告标题的日期}{业务指定的时区}
\clozerow{41}{机器的默认设置}{断言的期望值}{日志的显示格式}{导出判断函数}
\clozerow{42}{两个提交版本}{两套业务规则}{两种进程时区}{两个输入日期}
\clozerow{43}{周日 23:59}{周一 00:01}{周日 00:01}{周一 23:59}
\clozerow{44}{提前导出}{照常导出}{暂停测试}{不导出}
\clozerow{45}{机器时区已统一}{两条均符合规则}{整周已逐分验证}{旧函数已通过}

\newpage
\GaokaoSection{第二节}{（共 10 小题；每小题 1.5 分，满分 15 分）}

阅读下面的交接记录，在无提示词的空白处填入 1 个适当的虚词或副词，在有提示词的空白处结合语境为括号内词语增补 1 个表示状态或次序的语素。

\begin{center}
{\heiti\bfseries 订单导出修复交接记录}
\end{center}

14:05，订单导出报错，值班员停止向其余实例发布新版本。14:12，值班员~\blank{46}~已升级的两台实例~\blank{47}~新版本回退~\blank{48}~旧版本。14:20，从两台实例分别取回当前的可执行文件，逐一计算 \texttt{SHA-256}，均与旧版发布清单所列的值相符。两台的回滚记录至此收口。发布单上的“继续发布”一项此时标为\mbox{“\blank[3.2em]{49}”}（暂停）。本轮回滚~\blank{50}~涉及这两台，其余实例一直运行旧版本。

交班前，修复提交已进入主分支，合并单从“待合入”改为\mbox{“\blank[3.2em]{51}”}（合入）。交班时又核对了一句：“初测报错的那份脱敏样例，在预发布环境有没有用修复版重新测~\blank{52}~？”测了一次，结果已记在交接表“初测”右边的\mbox{“\blank[2.8em]{53}”}（测）栏里。生产发布仍待排期，请~\blank{54}~部署完成后再用线上订单核对。只有拿到这次线上核对的结果，\blank{55}~能填写“生产验证”一栏。

当前线上运行的是回滚后的旧版本，接班后请按旧版本排查新报错。

% ============================================================================
% 第四部分  写作（满分 40 分）
% ============================================================================
\vspace{2mm}
\GaokaoPart{第四部分}{写作（共两节，满分 40 分）}

\GaokaoSection{第一节}{（满分 15 分）}

\noindent 56.~周五 17:58，技术负责人张工给李华发来消息：“周末顺手把订单接口整理一下，旧的 \texttt{total} 字段直接删了，周一上线。”

李华请你帮忙回复。他说：“这周末我排不开。旧版移动端还在读 \texttt{total} 展示订单金额。周一我能跟他对方案，什么时候能上线现在还定不下来。”

请你作为助理给李华写一则回复。回复中须包含：

\begin{enumerate}[label=（\arabic*）, leftmargin=3em, labelsep=0pt, itemsep=1pt, topsep=2pt]
  \item 一条李华可以直接发给张工的消息，交代周末安排、直接删除字段的影响、周一可商谈的事项及上线时间的安排；
  \item 给李华的一句简短说明，交代如何发送或替换其中的措辞。
\end{enumerate}

\noindent 注意：（1）两层文字合计 120 字左右；（2）请在答题卡的相应位置作答。

\smallskip
\essaylines{6}

\newpage
\GaokaoSection{第二节}{（满分 25 分）}

\noindent 57.~阅读下面材料，根据其内容和所给段落开头语续写两段，使之构成一篇完整的技术手记。

\begin{quote}
\kaishu\parindent=2em
校刊排印前一晚，编辑发来一张预览截图。采访稿里“桥”字的位置多了几个问号形的替换符，后面的字都好好的。我从投稿系统下载原稿，那个字也好好的。编辑说她已经重打过一遍，预览还是老样子。我请她把那份会出错的原稿也留下。

预览服务每次最多读取 256 字节。为方便核对，我让它在日志里另记每次读到的原始字节。第一块的末尾是 \texttt{E6}，第二块开头是 \texttt{A1 A5}。原稿里“桥”的 \texttt{UTF-8} 编码正好是 \texttt{E6 A1 A5}，三个字节一个没少。本地测试一直把整篇稿子一次交给预览程序，跑了十次，十次都对。

编辑又试了一次。这回她只在标题前加了三个空格。“桥”能看清了，下一行的“街”却坏了。她问我，空格要不要留着。
\end{quote}

\noindent 注意：（1）续写字数应为 220 字左右；（2）请按如下格式在答题卡的相应位置作答。

\continuationpara{Paragraph 1:}{我请她先撤掉那三个空格，把原稿按日志里的位置分成两块。}{8}

\continuationpara{Paragraph 2:}{校样送印前，编辑又发来那份没有加空格的稿子。}{8}

\end{document}
